\documentclass[11pt,fontset=fandol]{ctexart}
\usepackage[a4paper,margin=23mm]{geometry}
\usepackage{amsmath,amssymb,mathtools,bm,mathrsfs,esint}
\usepackage{graphicx,longtable,booktabs,array,calc}
\usepackage[normalem]{ulem}
\usepackage{xcolor,hyperref}
\hypersetup{hidelinks}
\setcounter{secnumdepth}{-1}
\newcounter{none}
\providecommand{\tightlist}{\setlength{\itemsep}{0pt}\setlength{\parskip}{0pt}}
\providecommand{\passthrough}[1]{#1}
\usepackage{amssymb,mathrsfs,esint,bm,newunicodechar}
\allowdisplaybreaks
\setlength{\emergencystretch}{3em}
\providecommand{\llangle}{\langle\!\langle}
\providecommand{\rrangle}{\rangle\!\rangle}
\newunicodechar{∞}{\ensuremath{\infty}}
\newunicodechar{∈}{\ensuremath{\in}}
\newunicodechar{∑}{\ensuremath{\sum}}
\newunicodechar{∫}{\ensuremath{\int}}
\newunicodechar{≤}{\ensuremath{\leq}}
\newunicodechar{≥}{\ensuremath{\geq}}
\newunicodechar{≠}{\ensuremath{\neq}}
\newunicodechar{→}{\ensuremath{\to}}
\newunicodechar{⇒}{\ensuremath{\Rightarrow}}
\newunicodechar{α}{\ensuremath{\alpha}}
\newunicodechar{β}{\ensuremath{\beta}}
\newunicodechar{γ}{\ensuremath{\gamma}}
\newunicodechar{λ}{\ensuremath{\lambda}}
\newunicodechar{φ}{\ensuremath{\varphi}}
\newunicodechar{ε}{\ensuremath{\varepsilon}}
\newunicodechar{σ}{\ensuremath{\sigma}}
\newunicodechar{θ}{\ensuremath{\theta}}
\newunicodechar{δ}{\ensuremath{\delta}}
\newunicodechar{η}{\ensuremath{\eta}}
\newunicodechar{π}{\ensuremath{\pi}}
\newunicodechar{ξ}{\ensuremath{\xi}}
\newunicodechar{Σ}{\ensuremath{\Sigma}}
\newunicodechar{Δ}{\ensuremath{\Delta}}
\newunicodechar{∂}{\ensuremath{\partial}}
\newunicodechar{∀}{\ensuremath{\forall}}
\newunicodechar{∃}{\ensuremath{\exists}}
\newunicodechar{√}{\ensuremath{\surd}}
\newunicodechar{−}{\ensuremath{-}}
\newunicodechar{μ}{\ensuremath{\mu}}
\newunicodechar{ρ}{\ensuremath{\rho}}
\newunicodechar{τ}{\ensuremath{\tau}}
\newunicodechar{ω}{\ensuremath{\omega}}
\newunicodechar{Ω}{\ensuremath{\Omega}}
\newunicodechar{χ}{\ensuremath{\chi}}
\newunicodechar{ψ}{\ensuremath{\psi}}
\newunicodechar{Γ}{\ensuremath{\Gamma}}
\newunicodechar{Φ}{\ensuremath{\Phi}}
\newunicodechar{Λ}{\ensuremath{\Lambda}}
\newunicodechar{ℝ}{\ensuremath{\mathbb{R}}}
\newunicodechar{ℤ}{\ensuremath{\mathbb{Z}}}
\newunicodechar{ℕ}{\ensuremath{\mathbb{N}}}
\newunicodechar{△}{\ensuremath{\triangle}}
\newunicodechar{⊗}{\ensuremath{\otimes}}
\begin{document}
\section{2020年杭州电子科技大学601数学分析考研真题}\label{ux5e74ux676dux5ddeux7535ux5b50ux79d1ux6280ux5927ux5b66601ux6570ux5b66ux5206ux6790ux8003ux7814ux771fux9898}

来源：\texttt{HDU/601数学分析/2020/2020年杭州电子科技大学601数学分析考研真题.pdf}。原
PDF 页序：1, 2, 3。

\section{杭州电子科技大学}\label{ux676dux5ddeux7535ux5b50ux79d1ux6280ux5927ux5b66}

\section{2020
年攻读硕士学位研究生招生考试}\label{ux5e74ux653bux8bfbux7855ux58ebux5b66ux4f4dux7814ux7a76ux751fux62dbux751fux8003ux8bd5}

\section{《数学分析》试题}\label{ux6570ux5b66ux5206ux6790ux8bd5ux9898}

（试题共四大题，共3页，总分150分）

姓名 报考专业 准考证号

【所有答案必须写在答题纸上，做在试卷或草稿纸上无效！】

\subsection{一、填空题（本大题共8小题，每小题5分，共40分）}\label{ux4e00ux586bux7a7aux9898ux672cux5927ux9898ux51718ux5c0fux9898ux6bcfux5c0fux98985ux5206ux517140ux5206}

\begin{enumerate}
\def\labelenumi{\arabic{enumi}.}
\tightlist
\item
  \(\lim _ { n \to \infty } \left( { \frac { 1 } { n ^ { 2 } + n + 1 } } + { \frac { 2 } { n ^ { 2 } + n + 2 } } + { \frac { 3 } { n ^ { 2 } + n + 3 } } + \cdots + { \frac { n } { n ^ { 2 } + n + n } } \right) =\)
  （ ）.
\end{enumerate}

2.设 \(f ( x ) { = } \arctan x - { \frac { 1 } { 1 + a x ^ { 2 } } }\)
，且 \(f ^ { \prime \prime } ( 0 ) = 2\) ，则a=（ ）.

3.设y=f(x)由方程 \(y - \sin ( x y ) = e ^ { x ( 1 - y ) }\) 确定，则
\(\lim _ { n \to \infty } n \left[ f \left( { \frac { 1 } { n } } \right) - 1 \right] = (\)
）.

4.曲线 \(y = x ^ { 2 } + 2 \ln x\) 在其拐点处的切线方程为( ）.

5.已知
\(f ( x ) = \int _ { 1 } ^ { x } { \sqrt { 1 + t ^ { 4 } } } d t\) 则
\(\int _ { 0 } ^ { 1 } x ^ { 2 } f ( x ) d x =\) （ ）.

6.幂级数
\(\sum _ { n = 1 } ^ { \infty } ( - 1 ) ^ { n - 1 } n x ^ { n - 1 }\)
在区间(-1,1)内的和函数 S(x)=（ ）.

7.对区域
\(D { = } \left\{ ( x , y ) \vert x ^ { 2 } + y ^ { 2 } < 1 \right\}\)
内任意按段光滑曲线L，若曲线积分

\(\int _ { L } { \frac { x d x - a y d y } { x ^ { 2 } + y ^ { 2 } - 1 } }\)
与路径无关，则a=（ ）.

8.设
\(\Sigma = \left\{ ( x , y , z ) \left| x + y + z = 1 , x \geq 0 , y \geq 0 , z \geq 0 \right. \right\}\)
，则 \(\iint _ { \Sigma } y ^ { 2 } d S = (\) ）.

\subsection{二、计算题（本大题共5小题，每小题6分，共30分）}\label{ux4e8cux8ba1ux7b97ux9898ux672cux5927ux9898ux51715ux5c0fux9898ux6bcfux5c0fux98986ux5206ux517130ux5206}

1．求极限
\(\lim _ { x \to 0 } { \frac { \ln ( 1 + x ) ^ { \frac { 1 } { x } } - 1 } { \ln ( 1 + x ) } }\)

2．计算积分
\(\int _ { - 2 } ^ { 2 } \operatorname* { m i n } \left\{ { \frac { 1 } { \left| x \right| } } , x ^ { 2 } \right\} d x\)
,其中min表示取最小值.

3.设 \(z = z ( x , y )\) 为方程 \(f ( x - y , y - z , z - x ) = 0\)
所确定的隐函数，求
\(\frac { \partial z } { \partial x } , \frac { \partial z } { \partial y }\)
和dz.

4.求函数 \(f ( x ) = \arctan { \frac { 2 x } { 1 - x ^ { 2 } } }\)
在x=0处的幂级数展开式.

5.设区域
\(D { = } \left\{ \left( x , y \right) \bigm | x ^ { 2 } + y ^ { 2 } \leq a x \right\}\)
，计算二重积分

\[
I = \iint _ { D } \Bigl ( x ^ { 3 } y ^ { 2 0 1 9 } + y ^ { 3 } \sqrt { x ^ { 2 0 2 0 } + y ^ { 2 0 2 0 } } + x \sqrt { x ^ { 2 } + y ^ { 2 } } \Bigr ) d x d y .
\]

\subsection{三、综合题（本大题共3小题，每小题10分，共30分）}\label{ux4e09ux7efcux5408ux9898ux672cux5927ux9898ux51713ux5c0fux9898ux6bcfux5c0fux989810ux5206ux517130ux5206}

1.设
\(f ( x ) = { \left\{ \begin{array} { l l } { | x | ^ { \alpha } \sin { \frac { 1 } { | x | } } , 0 < | x | \leq 1 , } \\ { 0 , \qquad x = 0 , } \end{array} \right. }\)

(1）问α在什么条件下，f''(0)存在？并求此时f(0)的值.

(2）问α在什么条件下，f'(x)在{[}-1,1{]}上有界？

(3）问α在什么条件下，f''(x)在{[}-1,1{]}上连续？

2．求参数k取何值时，方程 \(k x + \frac { 1 } { x ^ { 2 } } = 1\) 在
\(x \in ( 0 , + \infty )\) 上有唯一解？

3.计算曲面积分
\(I = \iint _ { \Sigma } \frac { a x d y d z + ( z + a ) ^ { 2 } d x d y } { \sqrt { x ^ { 2 } + y ^ { 2 } + z ^ { 2 } } }\)
，其中Σ为下半球面\(z = - \sqrt { a ^ { 2 } - x ^ { 2 } - y ^ { 2 } }\)
的上侧，a\textgreater0为常数.

\subsection{四、证明题（本大题共5小题，每小题10分，共
50分)}\label{ux56dbux8bc1ux660eux9898ux672cux5927ux9898ux51715ux5c0fux9898ux6bcfux5c0fux989810ux5206ux5171-50ux5206}

1.设 \(\{ a _ { n } \}\) 是递增有界正数列.求证：级数
\(\sum _ { n = 1 } ^ { \infty } \left( 1 - { \frac { a _ { n } } { a _ { n + 1 } } } \right)\)
收敛.

2.设
\(f ( x ) = { \frac { x + 2 } { x + 1 } } \mathrm { s i n } { \frac { 1 } { x } } , a > 0\)
为常数.试证：f(x)在(0,a)内非一致连续，而在 {[}a,+∞)上一致连续.

3.设函数f(x)在区间{[}0,1{]}上连续，在开区间(0,1)内可导，且f(1)=0，证明：存在
\(\xi \in ( 0 , 1 )\) ，使得
\(f ( \xi ) + ( 1 - e ^ { - \xi } ) f ^ { \prime } ( \xi ) = 0\)

4.确定函数
\(g ( \alpha ) = \int _ { 0 } ^ { + \infty } { \frac { \ln ( 1 + x ^ { 3 } ) } { x ^ { \alpha } } } d x\)
的定义域，并证明g(α)在该定义域上连续.

5.设
\(f ( x , y ) = { \left\{ \begin{array} { l l } { { \underline { { x } } } ^ { 3 } + y ^ { 3 } } \\ { { \overline { { x } } } ^ { 2 } + y ^ { 2 } } \\ { 0 , } \end{array} \right. } ( x , y ) \neq ( 0 , 0 )\)

(1)证明f(x,y)在(0,0)处连续；(2）证明 \(f ( x , y )\) 在(0,0)处不可微.
\end{document}
